
\documentclass{article}
\usepackage{colortbl}
\usepackage{makecell}
\usepackage{multirow}
\usepackage{supertabular}

\begin{document}

\newcounter{utterance}

\centering \large Interaction Transcript for game `taboo', experiment `low\_en', episode 13 with reka{-}flash{-}3.1{-}t0.0.
\vspace{24pt}

{ \footnotesize  \setcounter{utterance}{1}
\setlength{\tabcolsep}{0pt}
\begin{supertabular}{c@{$\;$}|p{.15\linewidth}@{}p{.15\linewidth}p{.15\linewidth}p{.15\linewidth}p{.15\linewidth}p{.15\linewidth}}
    \# & $\;$A & \multicolumn{4}{c}{Game Master} & $\;\:$B\\
    \hline

    \theutterance \stepcounter{utterance}  
    & & \multicolumn{4}{p{0.6\linewidth}}{
        \cellcolor[rgb]{0.9,0.9,0.9}{
            \makecell[{{p{\linewidth}}}]{
                \texttt{\tiny{[P1$\langle$GM]}}
                \texttt{You are playing a collaborative word guessing game in which you have to describe a target word for another player to guess.} \\
\\ 
\texttt{Rules:} \\
\texttt{(a) You have to reply in the form: CLUE: <some text>. Guesses from the other player will start with GUESS.} \\
\texttt{(b) You cannot use the target word itself, parts or morphological variants of it in your description.} \\
\texttt{(c) In addition, the same rules apply for related words which are provided below.} \\
\\ 
\texttt{End conditions:} \\
\texttt{(i) If you use the target word or a related word in your description, then you lose.} \\
\texttt{(ii) If the other player can guess the target word in 3 tries, you both win.} \\
\\ 
\texttt{Let us start.} \\
\\ 
\texttt{This is the target word that you need to describe and that the other player needs to guess:} \\
\\ 
\texttt{wheat} \\
\\ 
\texttt{Related words are:} \\
\\ 
\texttt{{-} grain} \\
\texttt{{-} grass} \\
\texttt{{-} flour} \\
\\ 
\texttt{Important: You are under time pressure, give short descriptions that are to the point!} \\
            }
        }
    }
    & \\ \\

    \theutterance \stepcounter{utterance}  
    & \multicolumn{4}{p{0.6\linewidth}}{
        \cellcolor[rgb]{0.9,0.9,0.9}{
            \makecell[{{p{\linewidth}}}]{
                \texttt{\tiny{[P1$\rangle$GM]}}
                \texttt{<reasoning>} \\
\texttt{I need to come up with a clue for the word "wheat" without using the word itself, any parts of it, or related words like grain, grass, or flour. Let me think about what wheat is known for.} \\
\\ 
\texttt{Wheat is a type of grain, but that's a related word, so I can't mention that. It's used to make bread, pasta, etc., but flour is also a related word. So I can't use flour either. Grass is another related word, but wheat is a grass, so maybe I can't use that either. Wait, the related words are grain, grass, flour. So I can't use those. So I need to think of other characteristics.} \\
\\ 
\texttt{Wheat is a cereal crop. It's a staple food. It's golden in color, maybe? Or maybe the shape of the plant? Wheat stalks have those little grains attached. But grain is a related word. So maybe not. What else? It's used in baking, but flour is out. Maybe the process? Like milling? No, that's related to flour. } \\
\\ 
\texttt{Wait, maybe think about the color. Wheat is often associated with golden or yellowish color. But is that too vague? The other player might guess something else. Alternatively, think about the product that's made from wheat but not flour. Like bread, but that's made from flour. Maybe pasta? But again, that's made from wheat flour. Hmm. } \\
\\ 
\texttt{Alternatively, think about the plant itself. Wheat is a grass, but grass is a related word. So can't use that. Maybe the part of the plant? Like the ear of wheat? But the ear is part of the plant. Wait, the rules say I can't use parts or morphological variants. So "ear" might be a part. Maybe not. } \\
\\ 
\texttt{Wait, the target word is wheat. So I can't use "wheat" itself, but maybe other words that describe it. Let's think of synonyms. Cereal? But is cereal considered a related word? The related words given are grain, grass, flour. Cereal is a broader term, but maybe it's allowed. But maybe the user wants to avoid that. } \\
\\ 
\texttt{Alternatively, think about the uses. Wheat is used in making bread, but bread is a product. But the clue can't mention products made from it. Wait, the rules say you can't use the target word or} \\
            }
        }
    }
    & & \\ \\

    \theutterance \stepcounter{utterance}  
    & & & \multicolumn{2}{p{0.3\linewidth}}{
        \cellcolor[rgb]{0.9,0.9,0.9}{
            \makecell[{{p{\linewidth}}}]{
                \texttt{\tiny{[GM$|$GM]}}
                \texttt{abort game} \\
            }
        }
    }
    & & \\ \\

\end{supertabular}
}

\end{document}
